\documentclass[10pt,a4paper]{amsart}
\usepackage[margin=19mm]{geometry}
\usepackage{amsmath,amssymb,mathtools}
\usepackage[scr=rsfs,cal=euler]{mathalfa}
\usepackage{xfrac}
\usepackage[unicode,hidelinks,bookmarks=false]{hyperref}
\newcommand{\ie}{i.e.\ }
\newcommand{\cf}{cf.\ }
\newcommand{\resp}{respectively\ }
\newcommand{\Subsection}{\subsection}
\providecommand{\nobreakdash}{\nobreak\hbox{-}}
\setlength{\emergencystretch}{2em}
\multlinegap0pt
\usepackage{tikz}
\usetikzlibrary{cd}
\usepackage{mathtools}
\usepackage{amssymb}
\usepackage{tikz}
\newtheorem{theorem}{Theorem}
\newtheorem{remark}{Remark}
\title{Arithmetic Kodaira--Spencer Class and Frobenius Liftings via Frobenius--Witt Cotangent Complex}
\author{Kanau Shimada}
\date{Source: arXiv:2608.21772v1 (CC BY 4.0)}
\begin{document}
\maketitle

\noindent\emph{Source and licence.} Arithmetic Kodaira--Spencer Class and Frobenius Liftings via Frobenius--Witt Cotangent Complex. Version arXiv:2608.21772v1 (CC BY 4.0), distributed by arXiv under the Creative Commons Attribution 4.0 International licence, http://creativecommons.org/licenses/by/4.0/, as recorded in the arXiv licence metadata for this exact version and shown on the abstract page https://arxiv.org/abs/2608.21772v1. Full licence notice: \texttt{licenses/arxiv-2608.21772-CC-BY-4.0.txt}.\par\smallskip

\noindent\textit{Editorial scope.} Counted unit: the article's theorem thm3.9 (source0.tex lines 223--229). The article's complete original proof of it is delivered with it (source0.tex lines 230--264, one \texttt{proof} environment of 35 lines). No mathematics is written, repaired or re-derived: the statement and the proof are the article's own lines, in the article's own notation, and no retained line is substituted or re-flowed.  Retained uncounted as the source-local context the proof needs: lines 159--164; lines 190--192.  Nothing was omitted from any retained span, so the package carries no omission record.  No retained line contains a citation, so the wrapper carries no bibliography and no bibliography entry.  The retained lines contain the article's own diagram code, which the wrapper typesets with the standard tikz packages; no image file is delivered and the package declares no asset dependency.  Editorial formatting only, in addition: an AMS \texttt{amsart} preamble that restates the article's own declarations and macros used by the retained lines, namely the environments \texttt{theorem}, \texttt{remark} on separate counters, together with the standard packages those lines need.  The retained environments print in this document as Theorem 1, Remark 1 in order of appearance, whereas the article prints the same items with the numbering of its own full document.  The complete native source of the article as distributed in the arXiv e-print for this exact version is bundled unchanged as \texttt{sources/208275/source0.tex}.
\begin{theorem}\label{thm3.2}
    Let $X$ be a flat $\mathbb Z_{(p)}$-scheme. Write $X_0:=X\times _{\operatorname{Spec}\mathbb Z_{(p)}}\operatorname{Spec}\mathbb F_p$ and  $X_1:=X\times _{\operatorname{Spec}\mathbb Z_{(p)}}\operatorname{Spec}\mathbb Z/p^2$. Then the following assertions are equivalent.\begin{enumerate}
        \item The arithmetic Kodaira--Spencer class $\kappa_X$ of $X$ is zero.
        \item There is a Frobenius lift on $X_1$.
    \end{enumerate}
\end{theorem}

\begin{remark}\label{remA.1}
    In the situation of Theorem \ref{thm3.2}, the above proof implies that there is a natural bijection between the set of Frobenius lifts on $X_1$ and the set of retractions of the canonical map $\mathcal O_{X_0}\to F\mathbb L_X$.
\end{remark}

\begin{theorem}\label{thm3.9}
    Let $f:X\to Y$ be a morphism of flat $\mathbb Z_{(p)}$-schemes. Let $f_n:X_n\to Y_n$ denote the pullback of $f$ along the canonical morphism $\operatorname{Spec}\mathbb Z/p^{n+1}\to \operatorname{Spec}\mathbb Z$ for $n\ge0$. Assume that there is a Frobenius lift on $Y_1$ with corresponding retraction $r:F\mathbb L_Y\to \mathcal O_{Y_0}$ (Remark \ref{remA.1}). Then the following assertions are equivalent.
    \begin{enumerate}
        \item The composition\[F_{X_0}^*\mathbb L_{X_0/Y_0}\xrightarrow[]{\kappa_{X/Y}}f_0^*F\mathbb L_Y[1]\xrightarrow[]{f_0^*r}f_0^*\mathcal O_{Y_0}[1]\simeq \mathcal O_{X_0}[1]\] is zero,
        \item there is a Frobenius lift on $X_1$ which is compatible with the given one on $Y_1$.
    \end{enumerate}
\end{theorem}

\begin{proof}
$(1)\Rightarrow (2)$. Assume that the composition $f_0^*r\circ \kappa_{X/Y}$ is zero. By the universal property of pushout, we get the following commutative diagram:
\[
\begin{tikzcd}
  F^*_{X_0}\mathbb L_{X_0/Y_0} \lbrack -1 \rbrack \arrow[r, "\kappa_{X/Y}\lbrack -1\rbrack "] \arrow[d,] 
    & f^*_0F\mathbb L_{Y} \arrow[d, ] \arrow[ddr, "f^*_0r", bend left=20] \\
  0 \arrow[r, ] \arrow[drr, bend right=20] 
    & F\mathbb L_X \arrow[dr, "q"] \\
  & & \mathcal O_{X_0}.
\end{tikzcd}
\]
The upper triangle in the above diagram yields the following commutative diagram:
\[
\begin{tikzcd}
  % 1行目: (0,2)の対象
  \mathcal O_{X_0} \arrow[d] \arrow[ddr, "\mathrm{id}", bend left=20] & \\
  % 2行目: (0,1)の対象
  f^*_0F\mathbb L_{Y} \arrow[d] \arrow[dr, "f^*_0r"] & \\
  % 3行目: (0,0)の対象 と (1,0)の対象
  F\mathbb L_X \arrow[r, "q"] & \mathcal O_{X_0},
\end{tikzcd}
\]
which implies that the map $q$ is a retraction of the canonical map $\mathcal O_{X_0}\to F\mathbb L_X$ and the corresponding Frobenius lift on $X_1$ is compatible with the one on $Y_1$ corresponding to $r$.

$(2)\Rightarrow (1)$. Assume that there is a Frobenius lift on $X_1$ which is compatible with the given one on $Y_1$. Let $q$ be the corresponding retraction of the canonical map $ \mathcal O_{X_0}\to F\mathbb L_X$. Since the Frobenius lift under consideration on $X_1$ is compatible with the given one on $Y_1$, we have the following commutative diagram:

\[\begin{tikzcd}
    f_0^*F\mathbb L_Y \arrow[d] \arrow[dr, "f_0^*r"] &\\
    F\mathbb L_X \arrow[r, "q"] & \mathcal O_{X_0}.
\end{tikzcd}\]
Since the composition
\[ F^*_{X_0}\mathbb L_{X_0/Y_0} \lbrack -1 \rbrack \xrightarrow[]{\kappa_{X/Y}[-1]}f_0^*F\mathbb L_Y\to F\mathbb L_X\]
 is zero, the composition $f^*_0r\circ \kappa_{X/Y}$ is also zero.

\end{proof}

\end{document}
